\section{The exact minimum support-certificate rank}\label{sec:rank}
We compute the least value of $R(\cL)$ in \eqref{eq:rank-invariant} over
lifts of a ball built from a fixed family of symmetric cones.
Fix a finite dictionary $\mathscr K$ of simple symmetric cones. A lift
through $\mathscr K$ may use any finite number of copies of each member
and nonnegative rays; the dictionary is fixed, not the list of factors
of an individual lift. For a non-ray dictionary member
let $r$ and $a$ be its rank and Peirce constant. Put
\begin{equation}\label{eq:dictionary-capacity}
 B=\max_{K\in\mathscr K}a(K)(r(K)-1)>0,
 \qquad q_B(s)=\left\lceil\frac{s-1}{B}\right\rceil.
\end{equation}
We call $B$ the \emph{dictionary capacity}; it is the largest
\emph{primitive capacity} $a(r-1)$ of a dictionary member. It bounds
mixed curvature \emph{per unit of dual rank} and differs from the largest
mixed-curvature capacity $a\lfloor r^2/4\rfloor$ of one factor
(Theorem~\ref{thm:curvature}): for real PSD matrices of order four the
two capacities are three and four.

\begin{theorem}[Least worst-case certificate rank]\label{thm:rank-frontier}
Let $s\geq2$. Over all exact finite relative-Slater affine lifts of
$\B_2^s$ through $\mathscr K$,
\begin{equation}\label{eq:rank-frontier}
 \inf_{\cL}R(\cL)=q_B(s).
\end{equation}
For each fixed lift there is a dense open semialgebraic subset
$U\subset\Sph^{s-1}$, whose complement has surface measure zero, such
that
\[
 \sum_i\jrank Y_i\geq q_B(s)
 \quad\text{for every }v\in U\text{ and every }Y\in\cD_{\cL}(v).
\]
An explicit lift attains minimum certificate rank $q_B(s)$ at every
support except one pole, where its minimum rank is one. Consequently
\begin{equation}\label{eq:essential-frontier}
 \inf_{\cL}\essinf_{v\in\Sph^{s-1}}r_{\cL}(v)=q_B(s).
\end{equation}
\end{theorem}

\begin{proof}[Lower bound]
Use Lemma~\ref{lem:selection} to choose a primal lift $X(u)$ of $u$ and
a minimum-total-rank certificate $Y(v)$ of $1-\ip uv$. On a common dense
open regular subset of the sphere these are $C^1$ and satisfy
\[
 \sum_i\ip{X_i(u)}{Y_i(v)}=1-\ip uv.
\]
At a contact $u=v$, the negative mixed derivative of this identity is
the sphere metric, of rank $s-1$. With $p_i=\jrank X_i(u)$ and
$q_i=\jrank Y_i(u)$ there, Lemma~\ref{lem:peirce-channel} gives
\begin{equation}\label{eq:rank-lower-chain}
 s-1\leq\sum_i a_ip_iq_i
 \leq\sum_i a_i(r_i-1)q_i
 \leq B\sum_iq_i.
\end{equation}
The middle inequality uses $p_i\leq r_i-1$ whenever $q_i>0$;
when $q_i=0$ both terms vanish. Rays contribute zero curvature.
Reduced proper faces have no larger $a_i(r_i-1)$ than the original
dictionary factor. Since the selected certificate has minimum total rank,
\eqref{eq:rank-lower-chain} implies the asserted bound for \emph{all}
certificates at that support, and hence at every support in the regular
set. The same bound gives both lower bounds in
\eqref{eq:rank-frontier} and~\eqref{eq:essential-frontier}.
\end{proof}

\subsection{A uniform Peirce-perspective construction}
The matching upper bound uses one construction valid in every simple
Euclidean Jordan algebra. For $x\in V$ put $L_x y=x\circ y$ and
$P(x)=2L_x^2-L_{x^2}$, the quadratic representation.
For an idempotent $c$, $P(c)$ is orthogonal projection onto $V(c,1)$.
We use the standard facts that $P(x)K\subset K$ and that $P(x)$ is a cone
automorphism when $x$ is invertible. These hold in every Euclidean Jordan
algebra; see~\cite[Chapters~II--IV]{FK1994} and the Peirce and quadratic
representation discussion in~\cite[pp.~1554--1555]{GS2010}.
We also need a rank consequence. Approximate any $x$ by invertible elements $x_n$. Cone automorphisms preserve extreme
rays, so $P(x_n)c$ has rank one when $c$ is primitive. The closedness of
the set of nonnegative elements of rank at most one gives
\begin{equation}\label{eq:quadratic-rank}
 P(x)c\succeq0,\qquad \jrank(P(x)c)\leq1.
\end{equation}
This argument does not use associative matrix multiplication.

\begin{lemma}[Primitive Peirce perspective]\label{lem:perspective}
Let $V$ be simple of rank $r\geq2$, let $c$ be primitive, and set $d=e-c$.
The half-Peirce space $W=V(c,1/2)$ has dimension $a(r-1)$; give it the
norm induced by the trace inner product, $\norm{w}^2=\ip ww$. For $w\in W$,
\begin{align}
 X(\delta,w,z)&=\delta c+\sqrt2 w+zd\succeq0
   \quad\Longleftrightarrow\quad
   \delta,z\geq0,\quad\delta z\geq\norm w^2,
       \label{eq:perspective-positive}\\
 \det X(\delta,w,z)&=z^{r-2}(\delta z-\norm w^2).
       \label{eq:perspective-det}
\end{align}
For $g\in W$ and $A>0$ the element
\begin{equation}\label{eq:rank-one-completion}
 Y=A P\left(c-\frac{g}{\sqrt2 A}\right)c
   =Ac-\frac{g}{\sqrt2}+\frac{P(g)c}{2A}
\end{equation}
is positive of rank one. Its component in $V(d,1)$ has trace
$\norm g^2/(4A)$.
\end{lemma}
\begin{proof}
Peirce multiplication gives $w^2=tc+u$, where $u\in V(d,1)$ and
\[
 t=\ip{w^2}c=\ip w{w\circ c}=\frac{\norm w^2}{2}.
\]
The definition of $P$ and $c\circ w=w/2$ give $P(w)c=u$.
Its trace is $\tr u=\tr w^2-t=t$. By
\eqref{eq:quadratic-rank}, for $w\ne0$ we have $u=tk$ for a primitive
idempotent $k\in V(d,1)$. Thus $w^2=t(c+k)$. Spectral calculus shows
that $w$ lies in the Peirce-one subalgebra with identity $c+k$:
the supports of $w$ and $w^2$ coincide. Since $c\circ w=w/2$,
we have $k\circ w=w/2$. The span of $c,k,w$ is therefore a rank-two
Jordan subalgebra, orthogonal to $d-k$, and in this subalgebra the two
eigenvalues of $\delta c+\sqrt2 w+zk$ have sum $\delta+z$ and product
$\delta z-2t=\delta z-\norm w^2$.
The remaining $r-2$ eigenvalues are all $z$. This proves both displayed
claims for $w\ne0$. For $w=0$ they follow from the diagonal element
$\delta c+zd$.

Expanding the quadratic representation by its definition proves
\eqref{eq:rank-one-completion}. The first expression and
\eqref{eq:quadratic-rank} show positivity and rank at most one; its
$c$-coefficient $A>0$ makes the rank exactly one. The calculation above
with $w=g$ gives $\tr P(g)c=\norm g^2/2$, proving the trace formula.
\end{proof}

\begin{proof}[Attainment in Theorem~\ref{thm:rank-frontier}]
Choose a dictionary member attaining $B$, and use its notation from
Lemma~\ref{lem:perspective}. Put $q=q_B(s)$ and partition the first
$s-1$ real coordinates into $q$ nonempty groups, each of size at most $B$.
Isometrically embed each group into a subspace $W_G\subset W$.
Write a ball point as $x=(w_1,\ldots,w_q,b)$ and impose
\begin{equation}\label{eq:rank-attaining-lift}
 X_G=(1-b)c+\sqrt2 w_G+z_Gd\succeq0,
 \qquad \sum_{G=1}^q z_G=1+b.
\end{equation}
The variables $z_G$ are auxiliary and $w_G\in W_G$.
If the constraints hold, Lemma~\ref{lem:perspective} gives
$\sum_G\norm{w_G}^2\leq(1-b)\sum_Gz_G=1-b^2$, hence $x\in\B_2^s$.
Conversely, if $b<1$ and $x\in\B_2^s$, choose
$z_G\geq\norm{w_G}^2/(1-b)$ with sum $1+b$, which is possible because
$\sum_G\norm{w_G}^2/(1-b)\leq1+b$. At $b=1$ all $w_G=0$ and any
nonnegative $z_G$ of sum two works. Thus the projection is exact. At $x=0$ with $z_G=1/q$, all blocks
are positive definite, which proves Slater feasibility.

Let $v=(g_1,\ldots,g_q,\eta)\in\Sph^{s-1}$, with $g_G\in W_G$.
Every certificate admits the Peirce decomposition
\[
 Y_G=A_Gc+u_G+H_G,\quad u_G\in W,\quad H_G\in V(d,1),
\]
and positivity implies $A_G\geq0$ and $H_G\succeq0$.
Comparing coefficients on the whole affine slice
\eqref{eq:rank-attaining-lift} gives
\begin{equation}\label{eq:certificate-coefficients}
 \operatorname{proj}_{W_G}u_G=-g_G/\sqrt2,\quad
 \tr H_G=\gamma=\frac{1-\eta}{2},\quad
 \sum_G A_G=\frac{1+\eta}{2}.
\end{equation}
Indeed, varying the $z_G$ with fixed sum makes their trace coefficients
equal; the constant and $b$ coefficients give the two remaining scalar
identities. When $q=1$ these two comparisons give them directly.
Conversely, these conditions imply the whole-slice identity.

If $\eta<1$, define
\[
 A_G=\frac{\norm{g_G}^2}{2(1-\eta)}.
\]
For $g_G\ne0$, take the rank-one completion in
\eqref{eq:rank-one-completion}; its $V(d,1)$ component has trace
$\norm{g_G}^2/(4A_G)=\gamma$. For $g_G=0$, take
$Y_G=\gamma k_G$ with any primitive $k_G\in V(d,1)$.
The identity $\sum_G\norm{g_G}^2=1-\eta^2$ gives the prescribed sum of
the $A_G$. Thus every block has rank one and the total rank is $q$.
Conversely, $\gamma>0$ forces every block of every certificate to be
nonzero, so the minimum total rank is exactly $q$.
This includes the south pole $\eta=-1$.

At the north pole $\eta=1$, positivity and $\tr H_G=0$ give $H_G=0$.
Positivity then gives $u_G=0$: a positive element orthogonal to $d$
belongs to the face $\R_+c$. Concentrating $\sum_G A_G=1$ in one block
gives total rank one, which is minimal. If $W_G\ne W$, a
certificate may have extra half-Peirce components orthogonal to $W_G$.
Our construction sets them to zero, while the positive trace $\gamma$ of
the $V(d,1)$ component still forces each block to be nonzero. We do not
claim that certificates are unique. The exceptional pole has surface measure zero for
$s\geq2$, proving~\eqref{eq:essential-frontier} and finishing the theorem.
\end{proof}

\subsection{Examples and the order of quantifiers}
The primitive capacity $a(r-1)$ is $r-1$, $2(r-1)$, and $4(r-1)$ for
real, complex, and quaternionic Hermitian PSD matrices of order $r$, and
$m-2$ for a Lorentz cone $Q_m$. For the Albert cone $H_+^3(\OO)$ it is
$8(3-1)=16$. Consequently, lifts by any finite number of Albert factors
give the exact value $\lceil(s-1)/16\rceil$. The proof uses neither an associative
representation of the Albert algebra nor an embedding into a real PSD
cone.

In Theorem~\ref{thm:rank-frontier} the lift is fixed before the support
(the objective); reversing this order can change the value:
\begin{equation}\label{eq:quantifier-reversal}
 \inf_{\cL}\sup_v r_{\cL}(v)=q_B(s),\qquad
 \sup_v\inf_{\cL}r_{\cL}(v)=1.
\end{equation}
For the second identity rotate the construction so that the prescribed
$v$ is its north pole; the first is Theorem~\ref{thm:rank-frontier}.
The universal lower bound one follows from the nonzero support slack.
Likewise, the theorem minimizes over the full certificate fiber at each
support, not over certificates \emph{selected in advance} by an
algorithm. For $s=1$, whenever a non-ray factor is available, choose
orthogonal primitive idempotents $c,k$ and the section $X=tc+(1-t)k$,
$0\leq t\leq1$, with output $2t-1$. Its two support slacks have the
rank-one certificates $2k$ and $2c$, so the minimum is one; the formula
$\lceil(s-1)/B\rceil=0$ does not apply.
